\section{Attacks on Self-Writing Memory}
\label{sec:transfer}

This section describes the attacks. The first exploits a partition layout whose regions grow with
occupancy, and it motivates the security definition; the others exploit channels that
survive a layout built to close it. Experimental results for all of them are reported in
\cref{sec:eval}.

\subsection{Leakage of a Growing Partition Layout}

An IVF index is naturally stored with the records of each partition in a contiguous region,
because a search reads whole partitions and contiguous storage makes those reads efficient. When
regions grow as records are added, the snapshot at $\tau$ reveals the size of every partition and
the region that holds each sealed record. After $\tau$, the offsets of writes reveal which
partitions receive new records. Neither observation requires a key: encryption hides the
content of a record but not its location.

\subsection{Structure Transfer}

The attack consists of three steps.

\paragraph{Induction}
The adversary selects a topic $t$ from a known set and causes the agent to receive a task on $t$,
such as a document to summarize or a report to triage. The agent writes a record in its own words.
The adversary knows the topic and the approximate time of the induced record, but neither its
content nor its partition; unlike classical file injection \cite{zhang2016all}, it does not
choose the stored text. That supplied content causes a persistent memory write is documented
against deployed assistants \cite{rehberger2024spaiware,unit42memory} and across memory backends
\cite{das2026trojan}, and \cref{sec:eval-induction} measures it for a mem0-style pipeline
\cite{chhikara2025mem0} driven by local models.

\paragraph{Localization}
A flush of $w$ records reveals the set of partitions that received a record. The adversary repeats
the induction and combines the observed sets with one of two estimators.

\emph{Intersection.} If several induced records are known to fall into the same partition,
intersecting the partition sets of their flushes isolates that partition. If the other $w-1$
records of each flush fall into partitions uniformly at random, an unrelated partition appears in
a flush with probability $q = 1 - (1 - 1/B)^{w-1}$ and survives $r$ intersections with probability
$q^{r}$. The expected number of unrelated survivors, $(B-1)q^{r}$, falls below one once
$r > \ln(B-1) / \ln(1/q)$. The required number of repetitions grows with $w$, so batching raises
the cost of the attack without preventing it. An adversary generally cannot tell which induced
records share a partition, so this estimator is retained only as an upper bound.

\emph{Frequency analysis.} Without attribution, the adversary computes, for each topic $t$ and
partition $b$, the fraction of flushes carrying an induced record on $t$ in which $b$ receives a
write. A partition holding records on $t$ receives writes in the flushes of $t$ more often than in
the flushes of other topics. Each partition is labeled with the topic for which this rate is
highest, provided the rate exceeds the mean over all topics by twice the estimated sampling noise.
Padding writes to uniformly random partitions raise every rate by the same expected amount; they
add noise the adversary must overcome with a larger budget, but they do not remove the difference
between topics.

\paragraph{Transfer, and what it requires}
The adversary assigns each cohort-H record the label of the partition that holds it. By
\cref{lem:transfer} the expected accuracy of this assignment is the same on cohort H as on cohort
P, so the attack exposes the two cohorts equally, except to the extent that the agent's subject
matter has drifted between them or the adversary's taxonomy differs from the user's.

The step that makes this work is the lemma's third hypothesis, knowing which partition holds a
given cohort-H record, which a growing layout supplies in its snapshot and a sealed fixed-size
layout does not. The rest of this section is an inventory of channels that supply it anyway.

\subsection{Deployed Stores, and Time-Placed Controls}
\label{sec:transfer-general}

The growing-region layout is not particular to a research prototype: widely deployed vector stores
place records by content and persist that placement in files whose changed byte ranges a
size-preserving encryption layer such as \texttt{gocryptfs} or Cryptomator does not hide. FAISS
writes each inverted list contiguously, Chroma updates HNSW node and neighbor-list files
incrementally, and LanceDB stores IVF-PQ lists in a columnar dataset.

Placing by content is necessary for the attack but not sufficient, and the difference is what
\cref{sec:eval-real} measures. Transfer additionally needs the layout to disclose \emph{which}
partition holds a given stored record, at the granularity the adversary observes: a store that
rewrites its whole index on every save discloses nothing at block granularity, and the blocks a
graph-index insert touches depend on graph structure and insertion order as much as on content.
Stores that place by arrival time, such as an unindexed LanceDB table or a sqlite-vec table without
a partition key, are the cleanest test of the premise, and serve as negative controls.

\subsection{The Refusal Channel}
\label{sec:attack-refusal}

A store that bounds per-partition capacity must do something when a partition fills. If it refuses
the write and writes nothing, the refusal is observable, and by \cref{prop:capacity} what it
discloses is the partition's live load, which is dominated by cohort H once the adversary accounts
for its own induced contribution. The attack is direct: induce writes on a topic and watch whether
the region that topic's records occupy is rewritten. \Cref{sec:eval-refusal} measures it, and
measures the padded refusal that closes it.

\subsection{Derived Writes and Deduplication Membership Inference}
\label{sec:attack-derived}

\Cref{prop:derived} says a store whose operations depend on the content \emph{or the occurrence} of
what the agent does leaks cohort H to a feedback agent. The occurrence half admits a concrete
attack. Deduplication as a membership side channel is not itself new: Harnik, Pinkas, and
Shulman-Peleg showed that cross-user deduplication in cloud storage lets one user learn whether
another has already stored a file \cite{harnik2010side}. What is new here is the instantiation in
self-writing agent memory, where the adversary supplies the ``file'' as a document for the agent to
remember, the deduplication decision is made by the agent's own extraction pipeline, and the
disclosure is dated relative to a compromise.

Agent memory pipelines avoid storing a fact twice, and they do so by consulting what is already
stored. An adversary that supplies a document and can see whether a write happened therefore
learns whether a matching fact was already in memory. Because the adversary chooses the document
it chooses the fact, and a fact present before $\tau$ is a cohort-H membership disclosure.
Nothing is decrypted; the signal is the occurrence of a write. Against a store that flushes when
records are staged, the occurrence of a flush is the answer. Against one that flushes on a clock,
a flush happens either way and the channel is closed. \Cref{sec:eval-dedup} measures both sides.

The precondition holds beyond our prototype. Reading the shipped source of nine
open-source agent-memory frameworks, four make whether a write reaches storage depend on whether
the content is already stored: mem0 passes the most similar existing memories to its extraction
call and then drops any extracted fact whose hash is already present, so a known fact never
reaches the vector store; CrewAI asks a model for a consolidation plan whose \emph{insert} flag
decides whether a record is written; LangMem's memory manager patches existing memories and
creates only novel ones; and Graphiti deduplicates extracted entities against the graph before
writing them. The supplementary material lists each framework, its version and default store, and
the file and line at which the mechanism was found; a framework is recorded as vulnerable only
where the mechanism was located in its source.

\subsection{Multiple Versions}
\label{sec:attack-multisnapshot}

An adversary that retains several versions of the store file diffs consecutive versions and reads
off the regions rewritten in between. This is the write channel applied \emph{before} $\tau$, and
it is what supplies the cohort-H membership that \cref{sec:transfer} otherwise lacks: regions
labeled after $\tau$ then give those earlier writes their topics. The attack is parameterized by
the capture granularity $g$ of \cref{sec:granularity}, and \cref{sec:eval-multisnapshot} reports it
as a function of $g$, together with the session-boundary re-encryption that defeats it when the
adversary cannot capture more finely than a session.

\subsection{Counting Persisted Per-Keyword State}
\label{sec:attack-counting}

Forward-private schemes persist per-keyword state, such as Sophos's update counters
\cite{bost2016sophos} and Diana's counter-indexed GGM keys \cite{bost2017forward}, and in the local
setting that state is in the snapshot, so an adversary with auxiliary keyword frequencies can match
counters to keywords. Ties among the many rare keywords bound it: on D3, keyword recovery averages
0.05 and topic accuracy never exceeds the most-frequent-class baseline, so the channel that makes
\cref{prop:fp} a separation is weak in practice. The supplementary material gives the sweep.
